\doublespacing % Do not change - required

\chapter{Introduction}
\label{ch1}

%%%%%%%%%%%%%%%%%%%%%%%%%%%%%%%%%%%%%%%
% IMPORTANT
\begin{spacing}{1} %THESE FOUR
\minitoc % LINES MUST APPEAR IN
\end{spacing} % EVERY
\thesisspacing % CHAPTER
% COPY THEM IN ANY NEW CHAPTER
%%%%%%%%%%%%%%%%%%%%%%%%%%%%%%%%%%%%%%%


\section{Background}
\label{sec:background}

\ac{LLM}s have become a major component of modern natural-language processing systems, supporting applications including conversational interaction, text generation, summarisation, question answering and reasoning. However, their increasing capability has been accompanied by substantial growth in model size and inference cost. Although the computational expense of training large models is considerable, deployment introduces a distinct set of engineering constraints: model parameters must remain accessible during inference, intermediate activations must be processed throughout the network, and autoregressive generation requires previously computed attention states to be retained. Consequently, practical \ac{LLM} deployment is constrained not only by arithmetic throughput, but also by memory capacity, memory bandwidth and the efficiency with which model data are moved through the hardware.


\subsection{\ac{LLM} Deployment Challenges}
\label{subsec:llm-deployment}

The most persistent component of the inference-memory footprint is the model weight set. Transformer-based \ac{LLM}s contain large collections of parameters within the attention and feed-forward sub-networks, with linear transformations accounting for a substantial proportion of the model parameters. When these weights are stored using 16-bit floating-point formats such as FP16 or BF16, models containing billions of parameters require several gigabytes of memory before runtime activations or the \ac{KV} cache are considered.

A first-order indication of this deployment challenge can be obtained directly from the number of model parameters. If a model contains \(N\) parameters and each parameter is represented using \(b\) bits, the theoretical raw storage requirement of the weights can be approximated by Equation~\ref{eq:weights}:

\begin{equation}
S_{\mathrm{weights}}
=
N\frac{b}{8}
\quad \text{bytes}.
\label{eq:weights}
\end{equation}

For example, a three-billion-parameter model represented using 16-bit floating-point weights requires approximately

\[
3\times10^{9}\times\frac{16}{8}
\approx 6~\mathrm{GB}
\]

for the parameter payload alone. At the same numerical precision, a model containing 7 billion parameters requires approximately \(14~\mathrm{GB}\), a 70-billion-parameter model approximately \(140~\mathrm{GB}\), and a model at the 405-billion-parameter scale approximately \(810~\mathrm{GB}\). These estimates represent only the raw parameter payload and exclude runtime activations, the \ac{KV} cache, temporary computational buffers, framework overheads, memory fragmentation and other allocations required during inference. The actual deployment-memory requirement is therefore greater than that suggested by the model weights alone.

This scaling creates an immediate \emph{memory-capacity} problem. A model can only be executed entirely on a GPU if the tensors required during inference can be accommodated within the available device memory. Models that exceed the memory capacity of a single accelerator may therefore require model sharding across multiple GPUs, CPU or storage offloading, smaller batch sizes, shorter context lengths or other deployment compromises. Although these strategies enable larger models to run, they can increase system complexity and introduce additional communication or data-transfer overhead. Reducing the memory footprint of the model can therefore directly broaden the range of hardware on which an \ac{LLM} can be deployed.

Memory capacity is not the only constraint. Even when a model fits within GPU memory, its parameters must repeatedly be transferred from device memory to the computational units responsible for executing the Transformer layers. This introduces a separate \emph{memory-bandwidth} challenge. Whereas memory capacity determines whether the model can be accommodated on the device, memory bandwidth affects how rapidly the required model data can be supplied during inference.

The impact of memory traffic differs between the two principal phases of autoregressive inference. During prompt processing, commonly referred to as the \emph{prefill} phase, multiple input tokens can be processed in parallel and the resulting matrix operations typically exhibit comparatively high computational intensity. During autoregressive \emph{decode}, by contrast, one or a small number of new tokens are generally processed at each generation step. Although the Transformer layers must still access the model parameters for every generated token, less arithmetic work is performed per weight access. Consequently, the repeated movement of model weights can become an important contributor to decoding latency, particularly when inference is limited by memory bandwidth rather than raw floating-point computational throughput.

The inference-memory footprint also extends beyond the model weights. Intermediate \emph{activations} are generated dynamically as information propagates through the attention and feed-forward components of each Transformer block. Their memory requirements depend on factors including batch size, sequence length and hidden-state dimensionality. Unlike weights, which remain fixed after training and can therefore be quantized offline, activations depend on the current input and must be processed during runtime. Reducing their precision consequently requires a runtime quantization strategy and introduces an additional source of quantization error within the inference computation.

A further source of inference memory is the \ac{KV}. During autoregressive generation, the self-attention mechanism produces key and value representations for each processed token. Recomputing these representations for all preceding tokens whenever a new token is generated would introduce substantial redundant computation. Instead, previously computed key and value tensors are retained in memory and reused during subsequent decoding steps. This substantially improves computational efficiency, but creates an additional memory requirement that is separate from the static model-weight footprint.

Unlike model weights, whose storage requirement remains approximately constant once the model architecture and numerical precision are fixed, the size of the \ac{KV} cache increases with the number of tokens retained in the context. In simplified form,

\[
S_{\mathrm{KV}}
\propto
L,
\]

where \(S_{\mathrm{KV}}\) denotes the memory occupied by the cache and \(L\) denotes the retained sequence length. In practice, the total cache size also depends on the number of Transformer layers, the attention architecture, batch size and the numerical precision used to represent the cached key and value tensors. Consequently, longer contexts and larger batches can substantially increase the memory required during autoregressive generation.

This distinction is important because weight compression alone does not reduce every component of the inference-memory footprint. Weight-only quantization can substantially reduce the fixed storage occupied by model parameters, while leaving activations and the \ac{KV} cache at higher precision. As context length increases, the relative importance of KV-cache storage can therefore become greater. Reducing the precision of weights, activations and cached key--value tensors represents progressively broader levels of quantization scope, each offering additional opportunities for memory reduction while introducing additional sources of approximation error.

These deployment constraints remain relevant even for comparatively compact \ac{LLM}s. This dissertation therefore uses Llama-3.2-3B-Instruct as its principal evaluation model. At approximately three billion parameters, it provides a useful experimental scale: the model is sufficiently large for parameter storage, memory traffic, activation precision and KV-cache behaviour to remain meaningful engineering concerns, while remaining sufficiently compact to permit repeated calibration, quantization and evaluation within the computational resources available to this project.

The deployment problem considered in this dissertation should therefore not be viewed solely as a question of computational throughput. Practical \ac{LLM} inference is governed by the interaction between model size, memory capacity, memory bandwidth, dynamically generated activations and the growing \ac{KV} cache. Reducing the numerical representation cost of these tensors provides a natural route towards lowering deployment requirements, motivating the investigation of low-bit quantization.

\subsection{Motivation for low-bit quantization}
Quantization provides a direct approach to reducing the numerical representation cost of neural-network tensors. Rather than storing and processing every value using a high-precision floating-point representation, quantization maps values onto a smaller set of representable levels. In the context of LLM deployment, \ac{PTQ} is particularly attractive because it is applied to an already trained model and therefore does not require the original large-scale training process to be repeated. The theoretical storage benefit is straightforward. If a three-billion-parameter model requires approximately \(6~\mathrm{GB}\) for its weights at 16-bit precision, reducing the raw weight representation to 8 bits would ideally reduce this quantity to approximately \(3~\mathrm{GB}\), while a 4-bit representation would reduce it to approximately \[ 3\times10^{9}\times\frac{4}{8} \approx 1.5~\mathrm{GB}. \] Under this simplified calculation, moving from a 16-bit to a 4-bit weight representation corresponds to a fourfold reduction in the raw number of bits used to store the weights. Such reductions can be particularly valuable when deployment is constrained by GPU memory or by the amount of parameter data that must be moved during inference. 

The principal difficulty is that reducing numerical precision also reduces the number of values that can be represented. A quantizer must replace the original high-precision tensor values with nearby low-precision representations, introducing a quantization error. As bit width decreases, this approximation generally becomes more aggressive, and errors introduced into weights or intermediate tensors can propagate through subsequent Transformer layers. The practical objective is therefore not simply to minimise the nominal number of bits, but to achieve an appropriate balance between the reduction in representation cost and the resulting degradation in model quality. This motivates a central trade-off considered throughout this dissertation: \[ \boxed{ \mathrm{Model~Quality} \quad \longleftrightarrow \quad \mathrm{Compression} } \] However, bit width alone does not completely describe a quantized model. An equally important consideration is \emph{quantization scope}: which tensor classes within the inference process are actually represented at reduced precision. At the most conservative end of this design space, only the model weights are quantized. A configuration such as W4A16KV16 reduces weight precision to four bits while retaining the activations and KV cache at higher precision. Weight-only quantization can obtain substantial reductions in static parameter storage while limiting the number of locations at which quantization error enters the inference computation. 


\section{Research Aim}



\section{Project Objectives}



